\element{OPPH}{
  \begin{questionmult}{OPPH}\bareme{haut=2,p=0}
    Parmi les ondes ci-dessous, lesquelles sont des ondes planes progressives harmoniques ?
    \begin{multicols}{2}
      \begin{reponses}
        \bonne{$A\cos(\omega t-kx+\phi)$}
        \mauvaise{$A\cos(\omega t - k r)$}
        \mauvaise{$A\cos(\omega t)\cos(kx)$}
        \bonne{$\underline{A}e^{j(\omega t+\vv{k}.\vv{OM})}$}
      \end{reponses}
    \end{multicols}
  \end{questionmult}
}
\setgroupmode{OPPH}{cyclic}

\element{ondeSonore}{
  \begin{questionmult}{ondeSonore}\bareme{haut=2,p=0}
    Une onde sonore
    \begin{multicols}{2}
      \begin{reponses}
        \bonne{est une onde de surpression et de vitesse.}
        \mauvaise{est transverse.}
        \bonne{se propage forcément dans un milieu matériel.}
        \bonne{vérifier l'équation de d'Alembert dans l'approximation acoustique.}
      \end{reponses}
    \end{multicols}
  \end{questionmult}
}

\element{dAlembert}{
  \begin{question}{dAlembert1}\bareme{1}
    L'équation de d'Alembert s'écrit
    \begin{multicols}{2}
      \begin{reponses}
        \bonne{$\frac{\partial^2 u}{\partial x^2} - \frac{1}{c^2} \frac{\partial^2 u}{\partial t^2} = 0$}
        \mauvaise{$\frac{\partial^2 u}{\partial x^2} + \frac{1}{c^2} \frac{\partial^2 u}{\partial t^2} = 0$}
        \mauvaise{$\frac{\partial^2 u}{\partial x^2} - c^2 \frac{\partial^2 u}{\partial t^2} = 0$}
        \mauvaise{$\frac{\partial^2 u}{\partial x^2} + c^2 \frac{\partial^2 u}{\partial t^2} = 0$}
      \end{reponses}
    \end{multicols}
  \end{question}
}
\element{dAlembert}{
  \begin{question}{dAlembert2}\bareme{1}
    L'équation de d'Alembert s'écrit
    \begin{multicols}{2}
      \begin{reponses}
        \mauvaise{$\frac{\partial^2 u}{\partial t^2} - \frac{1}{c^2} \frac{\partial^2 u}{\partial x^2} = 0$}
        \mauvaise{$\frac{\partial^2 u}{\partial t^2} + \frac{1}{c^2} \frac{\partial^2 u}{\partial x^2} = 0$}
        \bonne{$\frac{\partial^2 u}{\partial t^2} - c^2 \frac{\partial^2 u}{\partial x^2} = 0$}
        \mauvaise{$\frac{\partial^2 u}{\partial t^2} + c^2 \frac{\partial^2 u}{\partial x^2} = 0$}
      \end{reponses}
    \end{multicols}
  \end{question}
}
\setgroupmode{dAlembert}{cyclic}

\element{vitessePhase}{
  \begin{questionmult}{vitessePhase}\bareme{haut=2,p=0}
    Pour une OPH se propageant dans le sens des x croissants, la vitesse de phase
    \begin{multicols}{2}
      \begin{reponses}
        \bonne{$=\frac{\omega}{k}$}
        \bonne{$=c$ pour une onde vérifiant l'équation de d'Alembert}
        \bonne{représente la vitesse à laquelle la phase se propopage.}
        \bonne{$=\lambda f$}
      \end{reponses}
    \end{multicols}
  \end{questionmult}
}

\element{impedanceAcoustique}{
  \begin{questionmult}{impedanceAcoustique}\bareme{haut=2,p=0}
    On supposer que $P_1$ et $v_1$ sont la surpression et la vitesse d'une onde plane progressive harmonique. L'impédance acoustique peut être définie comme
    \begin{multicols}{5}
      \begin{reponses}
        \bonne{$\mu_0 c$}
        \bonne{$\sqrt{\frac{\mu_0}{\chi_s}}$}
        \bonne{$\frac{P_1}{v_1}$}
      \end{reponses}
    \end{multicols}
  \end{questionmult}
}

\element{ondeEM}{
  \begin{questionmult}{ondeEM}\bareme{haut=2,p=0}
    Pour une onde électromagnétique
    \begin{multicols}{2}
      \begin{reponses}
        \mauvaise{$c=\frac{1}{\mu_0\epsilon_0}$}
        \bonne{$\vv{B}=\frac{\vv{k}\wedge\vv{E}}{\omega}$}
        \mauvaise{$\vv{\Pi}=\frac{\vv{E}\wedge\vv{B}}{\omega}$}
        \bonne{$\vv{E}$ et $\vv{B}$ sont orthogonaux.}
        \bonne{$\vv{E}$ et $\vv{k}$ sont orthogonaux.}
      \end{reponses}
    \end{multicols}
  \end{questionmult}
}

\element{hypotheses}{
  \begin{questionmult}{hypCorde}\bareme{haut=2,p=0}
    Quelles hypothèses sont nécessaires pour mettre en équation la corde vibrante
    \begin{multicols}{2}
      \begin{reponses}
        \mauvaise{le mouvement est uniquement horizontal}
        \bonne{les effets de la gravité sont négligés}
        \bonne{l'onde a une petite amplitude}
        \bonne{la corde est infiniment flexible}
        \bonne{les frottements sont négligés}
      \end{reponses}
    \end{multicols}
  \end{questionmult}
}
\element{hypotheses}{
  \begin{questionmult}{hypSonores}\bareme{haut=2,p=0}
    Quelles hypothèses sont nécessaires pour mettre en équation les ondes sonores
    \begin{multicols}{2}
      \begin{reponses}
        \mauvaise{les particules de fluides sont immobiles}
        \bonne{évolution isentropique}
        \bonne{l'onde a une petite amplitude}
        \bonne{le fluide est macroscopiquement au repos}
        \bonne{en l'absence de forces de viscosité}
      \end{reponses}
    \end{multicols}
  \end{questionmult}
}
\setgroupmode{hypotheses}{cyclic}